% basic nastaveni
\documentclass[12pt]{article}
\setlength{\parindent}{0pt}
\setlength{\parskip}{1em}
\usepackage[utf8]{inputenc}
\usepackage[czech]{babel}
\usepackage[T1]{fontenc}

% matematicke package
\usepackage{amsmath}
\usepackage{amsfonts}
\usepackage{amssymb}
\usepackage{geometry}
\geometry{margin=2.5cm}
\usepackage{pgfplots}
\pgfplotsset{campat=1.18}

\title{Maturitní otázka číslo 12: Posloupnosti a řady}
\author{}
\date{}

\begin{document}
\maketitle
\subsection*{Co to je posloupnost?}
Posloupnost je typem funkce, jejímž definičním oborem $D$ je buď množina $\mathbb{N}$ všech
přirozených čísel nebo jejím počátečním úsekem $\{1, 2, \dots, n\}$. Konečná posloupnost je
posloupnost s konečným definičním oborem, nekonečná posloupnost ho má nekonečný.

Konečná posloupnost s n-tým členem $a_n$ a s definičním oborem $\{1, 2, \dots, k\}$ se zapisuje
$\{a_1, a_2, \dots, a_k\}$ nebo $(a_n)_{n=1}^k$. Nekonečná posloupnost s n-tým členem $a_n$
se zapisuje $\{a_1, a_2, \dots, a_k, \dots\}$ nebo $(a_n)_{n=1}^\inf$.

\subsection*{Jak se zadávají posloupnosti?}
\subsubsection*{Výčet hodnot}
Tato metoda je vhodná pro konečné posloupnosti nebo pro ty, kde jde z prvních členů určit závislost,
ze které jde určit členy následující.

$(a_n)_{n=1}^6 = \{4, -3, 20, 22, -15, 42\}$, $(b_n)_{n=1}^\inf$ $= \{2, 4, 6, 8, \dots\}$

\begin{tabular}{|c|c|c|c|c|}
\hline
$n$   & 1 & 2 & 3  & 4 \\ \hline
$c_n$ & 3 & 9 & 27 & 81 \\ \hline
\end{tabular}
\subsubsection*{Graf}
Metoda grafem se využívá ve stejných případech jako metoda výčtem. Jeho výhoda oproti výčtu je jeho
názornost a lehčeji se z něj určuje závislost posloupnosti. Pro posloupnost
$(a_n)_{n=1}^7 = \{0, 1, 1, 2, 3, 5, 8\}$ by vypadal graf následovně:
\begin{center}
\begin{tikzpicture}
\begin{axis}[
    width=12cm,
    height=7cm,
    axis line = middle,
    xlabel = {$n$},
    ylabel = {$a_n$},
    ylabel style={
        at={(ticklabel* cs:0.5)},
        anchor=south,
        rotate=-90
    },
    xmin = 0, xmax = 8,
    ymin = 0, ymax = 10,
    xtick={0, 1, 2, 3, 4, 5, 6, 7, 8},
    ytick={0, 2, 4, 6, 8},
    grid=none,
    style={thick}
    ]
\addplot[
only marks,
mark=diamond*,
mark options={fill=red, draw=blue},
]
coordinates{
    (1,0) (2,1) (3,1) (4,2) (5, 3) (6, 5) (7, 8)
};
\end{axis}
\end{tikzpicture}
\end{center}
\subsubsection*{Vzorec pro n-tý člen}
Vzorec pro n-tý člen vyjadřuje vztah mezi hodnotou $n$ z definičního oboru a $a_n$ z oboru hodnot.
Po dosazení hodnoty $n$ do vzorce zjistíme hodnotu n-tého členu posloupnosti. Vzorec $a_n = 2n + 1$
nám dá posloupnost $\{3, 5, 7, \dots\}$.
\subsubsection*{Rekurentní vzorec}
Rekurentní vzorec určuje člen posloupnosti pomocí jednoho nebo více předcházejících členů. Součástí
každého rekurentního vzorce musí být také alespoň jeden člen posloupnosti, většinou je to to ten
první. Vzorec $a_n = a_{n-1} + 3, a_1 = 0$ nám dá posloupnost $\{0, 3, 6, \dots\}$.
\end{document}

